% Master Diagram 1 -- Dynamic CMOS gate (3-input NAND)
% Base for Q1 (pseudo-nMOS), Q2 (dynamic), Q3 (C2MOS), Q4 (domino), Q8 (n-p CMOS)
% NOTE: transistors are placed with "path ... node" (no backslash-n command)
% on purpose, so copy-paste tools cannot swallow characters. Valid LaTeX.
\documentclass[border=10pt]{standalone}
\usepackage{circuitikz}
\begin{document}
\begin{circuitikz}
% supply rail
\draw (-2,11) -- (2,11) node[right]{$V_{DD}$};
% precharge PMOS
\path (0,9.5) node[pmos] (P1){};
\draw (P1.S) -- (0,11);
\draw (P1.G) -- (-1.8,9.5) node[left]{$\phi$};
\draw (P1.D) -- (0,8.3) coordinate (Z);
\path (Z) node[circ]{};
\draw (Z) -- (1.8,8.3) node[right]{$Z$};
% n-block: three NMOS in series (A, B, C)
\path (0,6.8) node[nmos] (NA){}; \draw (NA.D) -- (Z);
\draw (NA.G) -- (-1.8,6.8) node[left]{$A$};
\path (0,5.0) node[nmos] (NB){}; \draw (NB.D) -- (NA.S);
\draw (NB.G) -- (-1.8,5.0) node[left]{$B$};
\path (0,3.2) node[nmos] (NC){}; \draw (NC.D) -- (NB.S);
\draw (NC.G) -- (-1.8,3.2) node[left]{$C$};
% clocked foot NMOS
\path (0,1.4) node[nmos] (NF){}; \draw (NF.D) -- (NC.S);
\draw (NF.G) -- (-1.8,1.4) node[left]{$\phi$};
\draw (NF.S) -- (0,0);
\path (0,0) node[ground]{};
\end{circuitikz}
\end{document}
